\subsection{Residual layers and effective depth}\label{sec:residual}

Consider the architecture
\begin{equation}\label{eq:residualarchitecture}
h_{\ell,i}=(1-\alpha)h_{\ell-1,i}
 +\alpha\tanh\!\left(b_{\ell,i}
       +\sum_jw_{\ell,ij}h_{\ell-1,j}\right),
\qquad 0\le\alpha\le1.
\end{equation}
The scalar $\alpha$ is known and belongs to the same
$2^{-b}\mathbb Z$ grid as the network parameters.
For a more finely encoded $\alpha$, replace $b$ by the larger
fractional bit length in the definition of $\bits$ and in the
preprocessing bound. Put $\tau=\alpha D$.
A skip step preserves the coordinate, so a geometrically distributed
jump can bypass many layers. This observation is essential when $D$
is large and $\alpha$ is small.

\begin{theorem}[Residual bound]\label{thm:residualupper}
Suppose that the actual maximum row norm is $S\le2$ and the biases
have absolute value at most one. The exact sampler can be extended
to \eqref{eq:residualarchitecture} with
$O(Dn^2\bits^6)$ preprocessing and $O(Dn^2\bits)$ stored bits.
Let $J$ count the root operation, the geometric-jump proposals, and
the input leaves in the compressed recursion. Then
\begin{align}
\E J&=O((1+\tau)^2) &&(S\le1),\label{eq:rescriticalupper}\\
\E J&=O((1+\tau)^2e^{7000(S-1)\tau}) &&(1<S\le2).\label{eq:ressuperupper}
\end{align}
With the scalar and geometric implementations described below,
the expected online bit work is $O(\bits^4\E J)$.
All implied constants are absolute; explicit choices are given in
the proof. The product bound
\eqref{eq:resproduct} is more informative for a particular network.
\end{theorem}

For fixed $1<s\le2$, the lower bound proved below has exponential
factor $(1+\alpha(s-1))^{2D}$. Since
\[
\frac{\alpha(s-1)}2\le\log(1+\alpha(s-1))
\le\alpha(s-1),
\]
the upper and lower exponential scales agree up to absolute
constant factors. At unit norm, the upper cost is polynomial in
$1+\alpha D$, and Theorem~\ref{thm:residualfused} lowers its exponent
to $35/32$. The critical exponent remains open for the full class.

\subsubsection{Intervals whose endpoints stay inside the output range}

Use $g=\max\{1,S\}$ and
$\varepsilon=2^{-P}$ with the precision $P$ in
\eqref{eq:precision}. In particular,
$\varepsilon\le[100(D+1)^4]^{-1}$.
Set $r_0=1$, $c_{0,i}=0$, and choose dyadic radii satisfying
\begin{equation}\label{eq:resrounded}
f_{\alpha,g}(r_{\ell-1})+2\alpha\varepsilon
\le r_\ell\le
f_{\alpha,g}(r_{\ell-1})+3\alpha\varepsilon,
\qquad
f_{\alpha,g}(r)=(1-\alpha)r+\alpha\tanh(gr).
\end{equation}
All stored radii and centers are rounded onto the fixed grid
$2^{-(P+b+6)}\mathbb Z$. Since a nonzero $\alpha$ is at least
$2^{-b}$, this grid is finer than $\alpha\varepsilon/64$.
Approximating the displayed scalar expression and then adding
$5\alpha\varepsilon/2$ supplies \eqref{eq:resrounded} on this
grid. Denominators therefore do not grow with the layer index.
When $\alpha=0$, the network is the identity and the sampler returns
$x_1$ immediately. Hence assume $\alpha>0$ below.

For a row, put $p=1-\alpha$ and retain the definitions
$\sigma,a,\rho,m,R$ from \eqref{eq:rowdata} and
\eqref{eq:midrad}. The residual interval has midpoint and radius
\[
M=p c_{\ell-1,i}+\alpha m,
\qquad
V=p r_{\ell-1}+\alpha R.
\]
Approximate $M$ to error at most $\alpha\varepsilon$, then project
that approximation onto $[-1+r_\ell,1-r_\ell]$ to obtain
$c_{\ell,i}$. The resulting invariant is
\begin{equation}\label{eq:resinvariant}
|h_{\ell,i}-c_{\ell,i}|\le r_\ell,
\qquad |c_{\ell,i}|+r_\ell\le1.
\end{equation}
Indeed, both endpoints $M\pm V$ lie in $[-1,1]$, because the
previous interval and both tanh endpoints do. Also
$V+\alpha\varepsilon\le r_\ell$.
Before projection the center belongs to $[M+V-r_\ell,M-V+r_\ell]$.
Projection onto $[-1+r_\ell,1-r_\ell]$ preserves membership in this
interval: its upper endpoint is at least $-1+r_\ell$, and its lower
endpoint is at most $1-r_\ell$.
This proves both inequalities in \eqref{eq:resinvariant}.
The projection may move the center by more than the approximation
error. The unweighted offspring proof below does not require a
small displacement from the midpoint.

The normalized residual sampler chooses the same-coordinate skip
with probability
\begin{equation}\label{eq:resskip}
p_\ell=\frac{p r_{\ell-1}}{r_\ell},
\end{equation}
and the centered tanh factory with probability
$\alpha R/r_\ell$. Its two constant masses complete the mixture
and contribute mean $(M-c_{\ell,i})/r_\ell$.
The interval invariant makes these masses nonnegative.
Its mean is therefore $(h_{\ell,i}-c_{\ell,i})/r_\ell$.
At the root, open this normalized sampler with probability $r_D$
and otherwise use constant masses $(1-r_D\pm c_{D,1})/2$.
The second invariant in \eqref{eq:resinvariant} makes this root
mixture valid.

\subsubsection{Compressing the skips}

Let $R_0=1$ and $R_{j+1}=f_{\alpha,g}(R_j)$.
This ideal sequence is decreasing. Concavity gives
$f'_{\alpha,g}(R_j)\le R_{j+1}/R_j\le1$, so the proof of
Lemma~\ref{lem:rounding} applies with per-step error
$3\alpha\varepsilon$:
\begin{equation}\label{eq:resroundingerror}
0< R_j\le r_j\le1,
\qquad 0\le r_j-R_j\le3j\alpha\varepsilon.
\end{equation}
The upper bound one follows from
$f_{\alpha,g}(1)+3\alpha\varepsilon<1$.
Moreover $\tanh(gr_j)\le r_j$: the ratio
$\tanh(gu)/u$ decreases with $u$, and this inequality holds at
$R_j$ because the ideal sequence decreases.

For $g=1$, Jensen's inequality applied to
$u\mapsto(1+u)^{-1/2}$ and the tanh lower bound gives
$f_{\alpha,1}(r)\ge r/\sqrt{1+\alpha r^2}$.
Monotonicity in $g$ consequently gives
\begin{equation}\label{eq:resradiuslower}
r_j\ge R_j\ge(1+\alpha j)^{-1/2}\qquad(g\ge1).
\end{equation}
Writing $r_\ell=f_{\alpha,g}(r_{\ell-1})+e_\ell$,
$0\le e_\ell\le3\alpha\varepsilon$, one obtains
\[
1-p_\ell
=\alpha+\frac{p(r_\ell-r_{\ell-1})}{r_\ell}
\le\alpha+\frac{3\alpha\varepsilon}{r_\ell}
\le2\alpha.
\]
Thus $h=\min\{2\alpha,1\}$ dominates every non-skip hazard.
From the current layer, draw a truncated geometric waiting time
with success probability $h$. Bypass the intervening layers.
At a proposal layer $\ell$, retain the proposal with probability
$(1-p_\ell)/h$. A rejected proposal takes one skip step; a retained
proposal uses the conditional non-skip mixture. If the waiting time
passes layer zero, probe the current input coordinate.
Independent geometric waiting times and thinning reproduce the
original skip probabilities exactly.

Truncated geometric sampling does not scan the skipped layers.
For a proposed gap $k\le D$, evaluate $(1-h)^k$ by fixed-point
binary powering and compare it with a lazy uniform random number.
Binary search locates the inverse distribution function.
There are at most $D+1$ boundaries. At precision
$m+3\lceil\log_2(D+2)\rceil$ the probability that any boundary
remains ambiguous is $O(2^{-m})$.
Binary powering needs $O(\log D)$ multiplications, and binary search
needs $O(\log D)$ comparisons. Schoolbook arithmetic with
$O(\bits+m)$ guarded bits therefore gives $O((\bits+m)^4)$ work per
refinement round and $O(\bits^4)$ expected work.
The arithmetic computes certified intervals and never decides an
equality of real numbers.

\subsubsection{A product bound for the compressed recursion}

For $0<r\le1$ and $1\le g\le2$, define
\begin{align*}
\chi_g(r)=\min\biggl\{&
\frac{\tanh(gr)}r
 \exp\!\left(\frac53g^2r^2+g^4r^4\right),\;
g(1+g^2r^2)e^{g^2r^2}\biggr\},\\
G_j&=1-\alpha+\alpha\chi_g(r_j).
\end{align*}
The first expression is Lemma~\ref{lem:offspring} before dividing
by the new radius. For the second expression, the full majority
has expected source count at most
$\rho(1+\rho^2)/\tanh\rho$, and the joint race factor is at most
$e^{\rho^2}$. Both bounds increase with the actual row norm,
so using $g$ is valid.
The expected number of normalized children at layer $j+1$ is at
most $(r_j/r_{j+1})G_j$.
After the root gate, the expected number of reached uncompressed
calls at layer $k$ is at most
$r_k\prod_{j=k}^{D-1}G_j$.
Each such call is a proposal with probability $h\le2\alpha$.
Nonnegative summation and the predictable trial-cost lemma give
\begin{equation}\label{eq:resproduct}
\E(J-1)\le
\prod_{j=0}^{D-1}G_j
+2\alpha\sum_{k=1}^{D}r_k\prod_{j=k}^{D-1}G_j.
\end{equation}
The empty product is one. This derivation also counts branches
that reach an input without any proposal.

\begin{lemma}[A radial comparison for the offspring factor]
\label{lem:resradial}
Put $d_g(r)=1-\tanh(gr)/r$. Whenever $0<r\le1$,
$1\le g\le2$, and $d_g(r)\ge0$,
\begin{align}
\chi_1(r)&\le1+4d_1(r)+5r^4,\label{eq:reschicritical}\\
\chi_g(r)&\le1+4d_g(r)+5r^4+6000(g-1).
\label{eq:reschisuper}
\end{align}
\end{lemma}
\paragraph{Proof idea.}
Relate the excess offspring cost to the amount by which one residual step
decreases its radius. Elementary Taylor bounds control the critical term and
its change when the row norm exceeds one.

\begin{proof}
Since $r-\tanh r=\int_0^r\tanh^2u\,du$ and
$\tanh^2u\ge u^2/(1+u^2)\ge u^2-u^4$,
\[
d_1(r)\ge r^2/3-r^4/5.
\]
Write $t=r^2$. All Taylor coefficients of
$E(t)=\exp(5t/3+t^2)$ are nonnegative. On $0\le t\le2/5$,
they give $E(t)\le1+5t/3+4t^2$, because
$E(2/5)=e^{62/75}<173/75$.
Here is one rational verification of this numerical inequality:
$e<68/25$ by the exponential series and its geometric tail,
$e^{1/6}\ge85/72$, and therefore
$e^{62/75}<e^{5/6}<4896/2125<173/75$.
Multiplication by $1-d_1(r)$ and
$d_1(r)\ge t/3-t^2/5$ now gives
$\chi_1(r)\le1+4d_1(r)+5t^2$.
For $2/5\le t\le1$, the second expression defining $\chi_1$
and positivity of the series for $(1+t)e^t$ give
$\chi_1(r)\le1+2t+(5/2)t^2$.
This is at most $1+4d_1(r)+5t^2$, since
$(17/10)t^2-(2/3)t\ge0$ for $t\ge2/5$.

Let $\delta=g-1$. If $\delta\le1/10$, differentiating the first
expression defining $\chi_g$ with respect to $g$ bounds its
derivative by 400. Indeed, $gr\le11/10$,
$5g^2r^2/3+g^4r^4<7/2$, and the remaining derivative factor is
less than 11; $e^{7/2}<36$ suffices.
The derivative of the second expression is
\[
e^{g^2r^2}(1+5g^2r^2+2g^4r^4)<40
\]
on the same interval. Both functions are therefore
400-Lipschitz in $g$, and so is their minimum.
Also $d_1(r)\le d_g(r)+\delta$, by the unit Lipschitz bound for
tanh. This proves \eqref{eq:reschisuper} with 404 in place of 6000
in this case. If $\delta\ge1/10$, the second expression gives
$\chi_g(r)\le10e^4<600$, which proves the stated bound.
\end{proof}

\paragraph{Proof idea.}
The residual radius supplies a telescoping potential for the offspring
product. Summing the remaining fourth-order and rounding terms over the
effective depth gives the stated bound for compressed skips.

\begin{proof}[Proof of Theorem~\ref{thm:residualupper}]
We prove the explicit prefactors $9e^{21}$ in the critical case
and $10e^{30}$ in the supercritical case.
Write $d_j=d_g(r_j)$ and
$F_j=f_{\alpha,g}(r_j)=r_j(1-\alpha d_j)$.
The preceding lemma, $\log(1+u)\le u$, and
$-\log(1-v)\ge v$ give
\begin{equation}\label{eq:reslogproduct}
\log G_j\le
4\log\frac{r_j}{r_{j+1}}+5\alpha r_j^4
+6000\alpha(g-1)+\frac{12\alpha\varepsilon}{F_j}.
\end{equation}
Here the last term accounts for
$r_{j+1}=F_j+e_{j+1}$ and
$4\log(1+e_{j+1}/F_j)\le12\alpha\varepsilon/F_j$.
By \eqref{eq:resradiuslower}, its sum over all layers is less than
one.

At $g=1$, Lemma~\ref{lem:radius} and convexity imply
\[
\frac{\tanh r}{r}\le(1+2r^2/3)^{-1/2}\le1-r^2/5,
\qquad 0\le r\le1.
\]
The last inequality follows from the chord between zero and one
and $\sqrt{3/5}<4/5$.
Thus the ideal critical recurrence obeys
$R_j\le(1+2\alpha j/5)^{-1/2}$.
Integral bounds and \eqref{eq:resroundingerror} yield
\begin{equation}\label{eq:resintegrals}
\alpha\sum_{j=0}^{D-1}r_j^4\le4,
\qquad
\alpha\sum_{j=1}^{D}r_j^5\le3.
\end{equation}
For clarity, the ideal sums are at most
$1+\int_0^\infty(1+2t/5)^{-2}\,dt=7/2$ and
$1+\int_0^\infty(1+2t/5)^{-5/2}\,dt=8/3$.
The rounding contributions are at most
$6\varepsilon D^2$ and $15\varepsilon D(D+1)/2$, respectively.
Summing \eqref{eq:reslogproduct} therefore gives
\[
\prod_{j=k}^{D-1}G_j\le e^{21}(r_k/r_D)^4.
\]
Insert this inequality into \eqref{eq:resproduct}, use
\eqref{eq:resintegrals}, and finally use
$r_D^{-4}\le(1+\tau)^2$. Adding the root operation proves
\eqref{eq:rescriticalupper}.

For $g=1+\delta$, the unit Lipschitz bound gives
$d_g(r)\ge r^2/5-\delta$. Since $0\le d_j\le1$,
\[
\alpha r_j^2d_j\le r_j^2-F_j^2
\le r_j^2-r_{j+1}^2+2e_{j+1}.
\]
Consequently
\[
\alpha\sum_{j=k}^{D-1}r_j^4
\le5+5\delta\alpha(D-k)+30\alpha(D-k)\varepsilon.
\]
Equation \eqref{eq:reslogproduct} now gives
\[
\prod_{j=k}^{D-1}G_j
\le e^{30}(r_k/r_D)^4
       e^{6025\delta\alpha(D-k)}.
\]
Put $a=1+\alpha\delta$. Concavity and
$\tanh(gr)\le g\tanh r$ show
$f_{\alpha,g}(r)\le a f_{\alpha,1}(r)$.
Induction gives $R_j^{(g)}\le a^jR_j^{(1)}$.
Hence, using the same rounding estimates,
\[
\alpha\sum_{j=1}^D r_j^5\le4e^{5\delta\tau}.
\]
Substitution in \eqref{eq:resproduct} proves
\eqref{eq:ressuperupper}, with room in the displayed constants.

The normalized construction is exact by induction on the depth.
Geometric thinning preserves its transition probabilities.
Almost sure termination follows first for the finite-depth
recursion and then for each scalar comparison; no expected-work
bound is assumed in this argument.
Each reached proposal has conditional expected scalar cost
$O(\bits^4)$ for $g\le2$. The hazard ratio introduces no small-$\alpha$
division in the cost: $\tanh(gr)/r\ge\tanh1>3/4$ implies
$1-p_\ell\ge\alpha/2$ under the chosen rounding precision.
The predictable trial-cost lemma charges scalar operations,
weighted index selection, and geometric jumps by $O(\bits^4\E J)$.
Centers, radii, and row data use $O(\bits)$ bits; their preprocessing
fits the stated budget.
\end{proof}

\subsubsection{A simple construction for large row norms}

The preceding estimates use small uncertainty radii to exploit
near-criticality. A simpler construction gives a useful branch
bound for every finite norm $s$.

\begin{proposition}[Quadratic single-row branch cost]
\label{prop:largerow}
For a row of norm $\sigma$ and bias $b_0$, set
$u=\sigma+|b_0|$. There is an exact tanh-row sampler using at
most $2\sigma(1+\sigma)$ expected recursive child requests.
If every residual row has norm at most $s$, put
$C_s=2s(1+s)$ and $b_s=1-\alpha+\alpha C_s$.
There is an exact residual sampler whose expected number of
input leaves plus active residual rows is at most
\begin{equation}\label{eq:largeresidual}
b_s^D+\alpha\sum_{j=0}^{D-1}b_s^j.
\end{equation}
In particular, for $s\ge1$ this is at most $4b_s^D/3$.
For each fixed $s$, expected bit work is
$O_s(\bits^4[1+b_s^D+\alpha\sum_{j<D}b_s^j])$.
For arbitrary $s\ge1$, Corollary~\ref{cor:largeresidualsharpbits}
gives the uniform bit bound $O(\bits_s^4b_s^D)$, where
$\bits_s=\bits+\lceil\log_2(s+2)\rceil$.
Its preprocessing consists of rational row sums and prefix tables
and takes $O(Dn^2\bits_s^2)$ bit operations.
\end{proposition}
\paragraph{Proof idea.}
At each active update, sample the bias and row through one scalar source.
Consecutive skip choices can be generated in a single geometric jump, so the
accounting counts active updates only.

\begin{proof}
The row sampler is Corollary~\ref{cor:comptanhrow} with $R=1$, the
children being the source signs. Its proof bounds the expected child
requests by $2\sigma\max\{1,u\}\le2\sigma(1+\sigma)$, and its
conditional charging covers a majority that stops early.

At each residual row, first select the nonlinear branch with
probability $\alpha$. Otherwise skip to the preceding layer.
Geometric waiting times implement these skips directly.
The expected number of descendant input leaves is at most $b_s^D$,
and the expected number of active rows is at most
$\alpha\sum_{j<D}b_s^j$. This proves \eqref{eq:largeresidual}.
For $s\ge1$, $C_s\ge4$ and the geometric sum contributes at most
$(b_s^D-1)/(C_s-1)\le b_s^D/3$.
Theorem~\ref{thm:quadraticradiusbits} gives scalar cost
$O((1+s)^2\bits_s^4)$ at each active row, excluding child
computations. Compressed skips and input leaves cost
$O(\bits_s^4)$ each. For fixed $s$, the predictable trial-cost
lemma proves the stated bit bound. Corollary~\ref{cor:largeresidualsharpbits}
absorbs the quadratic row factor into the branching sum.
\end{proof}



\begin{proposition}[The quadratic norm scale is necessary in one layer]
\label{prop:residuallargeslower}
Let $s\ge2$ and $0\le\alpha\le1$.
For an unknown source sign of mean $z$, the optimal worst-mean
source count for the target
$f(z)=(1-\alpha)z+\alpha\tanh(sz)$ satisfies
\[
\frac3{64}(1+\alpha s^2)\le U\le2(1+\alpha s^2).
\]
For one residual layer with zero biases, uniform row weights
$s/m$, and even $m\ge2$, the optimal worst-input expected
number of distinct probes satisfies
\begin{equation}\label{eq:residuallargeslower}
\frac1{130}(1+\alpha\min\{s^2,m\})
\le Q^*\le2(1+\alpha\min\{s^2,m\}).
\end{equation}
The uniform weights are assumed legal on the parameter grid.
Thus the one-layer scale $\alpha s^2$ is necessary as well as
sufficient, until the finite input size becomes the limiting
resource. This statement does not assert that the factor
multiplies optimally across arbitrary depth.
\end{proposition}
\paragraph{Proof idea.}
A skip-or-update choice gives the upper bound. For the lower bound, a mean
network retains a large derivative at zero; the transcript inequality
converts this derivative into a probe requirement.

\begin{proof}
The upper bound chooses the skip with probability $1-\alpha$
and the global tanh majority with probability $\alpha$.
Its source count is at most $1-\alpha+2\alpha s^2$.
For a finite input, one may instead scan all $m$ coordinates on
the nonlinear branch. Choosing the cheaper construction gives
$1-\alpha+\alpha\min\{2s^2,m\}$ expected probes and proves
the finite upper bound.
Endpoint coupling gives both lower-bound quantities at least
$(1-\alpha)+\alpha\tanh s\ge1/2$.
For the unknown-source problem, the second transcript inequality,
justified by the truncation argument preceding
\eqref{eq:unknownlower}, applies at $z=1/s$.
Since $\tanh1\ge1/2$ and $\sech^2 1\ge1/4$,
\[
U\ge\frac{1-s^{-2}}2\,2\alpha s^2\tanh1\sech^21
\ge3\alpha s^2/32.
\]
Combining this with $U\ge1/2$ proves the stated bound.

For the finite network, under independent inputs of mean $t$,
put $G(t)=\E_t\tanh(sM)$, $R=\tanh s$, and
$q_0=\min\{s^2,m\}$. The averaged output is
$H(t)=(1-\alpha)t+\alpha G(t)$.
The same weighted Jensen argument as in
\eqref{eq:derivativelower} gives
\[
h:=G'(0)\ge\frac{s}{\sqrt{1+3s^2/m}},
\qquad h^2\ge q_0/4.
\]
If $q_0\ge64$, set $T=2R/h\le1/2$.
Since $G(0)=0$, $G'(0)=h$, and $G(T)\le R$, the bending argument
before \eqref{eq:bending} gives some $t\in[0,T]$ with
$|G''(t)|\ge h^2/(2R)$.
The linear skip contributes no second derivative, so the finite
transcript inequality yields
\[
Q^*\ge\frac{3\alpha h^2}{16R}
\ge\frac{3\alpha q_0}{64}.
\]
Together with $Q^*\ge1/2$, this is stronger than the lower bound
in \eqref{eq:residuallargeslower}.
If $q_0<64$, the endpoint bound alone gives
$Q^*\ge1/2\ge(1+\alpha q_0)/130$.
\end{proof}

\begin{corollary}[Residual bit bound at arbitrary norm]
\label{cor:largeresidualsharpbits}
Let $s\ge1$, with $C_s$, $b_s$, and $\bits_s$ as in
Proposition~\ref{prop:largerow}.
The raw residual sampler has expected online bit work
\begin{equation}\label{eq:largeresidualsharpbits}
 O(\bits_s^4 b_s^D).
\end{equation}
Its preprocessing remains rational row sums and prefix tables,
at cost $O(Dn^2\bits_s^2)$.
In particular, $\alpha s^2D=O(\log\log n)$ and
$\bits_s=\log^{O(1)}n$ suffice for polylogarithmic expected
online bit work. This permits $s$ itself to exceed every fixed
power of $\log n$. To fit the original
$O(Dn^2\log^6n)$ preprocessing allowance, it is sufficient
that $\bits_s=O(\log^3n)$.
\end{corollary}
\paragraph{Proof idea.}
Use the sharp large-radius scalar routine at each active residual update.
The same compressed-skip accounting now charges its quadratic radius cost
together with all scalar precision work.

\begin{proof}
An active row has $u=|b_0|+\sigma\le1+s$ and therefore
costs $O((1+s)^2\bits_s^4)$ by
Theorem~\ref{thm:quadraticradiusbits}.
Its expected number of child requests remains at most $C_s$.
The expected number of leaves is at most $b_s^D$, and the
expected active-row count is at most
$\alpha\sum_{j<D}b_s^j$.
Compressed geometric skips have the same logarithmic bit cost
per visited branch. The stopping-cost lemma gives
\[
 O\!\left(\bits_s^4\left[
 1+b_s^D+(1+s)^2\alpha\sum_{j<D}b_s^j\right]\right).
\]
For $s\ge1$,
\[
 \frac{(1+s)^2}{C_s-1}\le\frac43,
 \qquad
 (1+s)^2\alpha\sum_{j<D}b_s^j
 \le\frac43(b_s^D-1),
\]
including $\alpha=0$ by continuity or direct inspection.
This proves \eqref{eq:largeresidualsharpbits}.
Finally $C_s-1\le4s^2$ and
$b_s^D\le\exp(4\alpha s^2D)$.
\end{proof}

\begin{corollary}[One residual layer]
\label{cor:largeresidualonelayerbits}
For any signed, biased row of width $n$, row norm at most
$s\ge1$, and $|b_0|\le1$, one residual output admits an
exact sampler using
\[
 O\!\left(\bits_s^4
 [1+\alpha\min\{s^2,n\}]\right)
\]
expected online bit operations after the same preprocessing.
For the legal uniform zero-bias rows of
Proposition~\ref{prop:residuallargeslower}, with $s\ge2$ and
even width, its input-query dependence is optimal up to an
absolute factor.
\end{corollary}
\paragraph{Proof idea.}
There is only one active nonlinear row. Mix its exact sample with the
directly accessible skip coordinate, and compare this algorithm with reading
the entire row.

\begin{proof}
Choose the identity input on the skip branch. On the nonlinear
branch, put $H=\max\{1,\sigma+|b_0|\}$ and use the
quadratic-radius sampler if $H^2\le n$; otherwise probe and
sum the entire rational row and sample
its tanh output with a certified comparison. The latter costs
$O(n\bits_s^4)$ in expectation. Since $H\le2s$, the chosen
cost is bounded by $O(\bits_s^4\min\{s^2,n\})$.
The query lower bound is
\eqref{eq:residuallargeslower}.
\end{proof}


\subsubsection{A matching law for residual mean networks}

Let $m\ge2$ be even. Use zero biases and put $w_{\ell,ij}=s/m$
on the same $m$ relevant coordinates in every row and layer.
The displayed weights are assumed to be legal on the dyadic grid.
For class lower bounds at fixed $s>1$, a downward dyadic
approximation to $s$ makes them legal and preserves the stated
exponential order.
Define
\begin{equation}\label{eq:reschainparameters}
\begin{aligned}
p&=1-\alpha,\qquad a=p+\alpha s,\qquad \lambda=a^D,\\
F_0(z)&=z,\qquad
F_{j+1}(z)=pF_j(z)+\alpha\tanh(sF_j(z)),\\
A&=\frac{\alpha s^3}{a}\sum_{j=0}^{D-1}a^{2j}.
\end{aligned}
\end{equation}
Here $0<s\le2$, so $a>0$; $\alpha=0$ is permitted.
Writing $M=m^{-1}\sum_i x_i$, direct induction gives
\begin{equation}\label{eq:resmeanidentity}
h_{D,1}=F_D(M)+p^D(x_1-M).
\end{equation}

\begin{theorem}[Uniform residual mean-network law]
\label{thm:resmeanlaw}
For the network \eqref{eq:resmeanidentity}, define
\[
\Xi=\lambda\min\left\{\sqrt{1+A},\frac m{\sqrt{1+A}}\right\}.
\]
Its optimal worst-input expected number of distinct input probes
satisfies
\begin{equation}\label{eq:resmeanlaw}
\frac\Xi{512}\le Q^*\le3\Xi.
\end{equation}
The constants are uniform in $s,\alpha,D,m$ in the stated range.
For an unknown source coin and target mean $F_D(z)$, the optimal
worst-mean source count $U$ satisfies
\begin{equation}\label{eq:resunknownlaw}
\frac1{81}\lambda\sqrt{1+A}\le U
\le2\lambda\sqrt{1+A}.
\end{equation}
These statements concern source or input queries; auxiliary
arithmetic is excluded from this particular measure.
\end{theorem}

\begin{lemma}[Residual range comparison]\label{lem:resrange}
For $0\le z\le1$,
\begin{equation}\label{eq:resrangesandwich}
\frac{\lambda z}{\sqrt{1+Az^2}}
\le F_D(z)\le
\frac{\lambda z}{\sqrt{1+Az^2/5}}.
\end{equation}
\end{lemma}
\paragraph{Proof idea.}
Compare one residual step with reciprocal-square maps above and below it.
Composing these inequalities gives explicit bounds for the range of the full
scalar chain.

\begin{proof}
Set $v=\alpha s/a$ and $\theta=\alpha s^3/a$.
The one-step ratio is
\[
\frac{p z+\alpha\tanh(sz)}{az}
=(1-v)+v\frac{\tanh(sz)}{sz}.
\]
The tanh lower bound and convexity of $(1+u)^{-1/2}$ give the lower
bound $(1+\theta z^2)^{-1/2}$ for this ratio.
For $0\le y\le2$, the coth inequality \eqref{eq:cothradius} and a
chord bound give
\[
\frac{\tanh y}y\le(1+2y^2/3)^{-1/2}\le1-y^2/10.
\]
The last chord uses $(1+8/3)^{-1/2}<3/5$.
Since $1-w/2\le(1+w)^{-1/2}$ for $w\ge0$, the ratio is at most
$(1+\theta z^2/5)^{-1/2}$.
Every iterate remains in $[0,1]$.
Iterating these reciprocal-square bounds gives
\eqref{eq:resrangesandwich}; the accumulated sum is precisely $A$.
\end{proof}

\begin{lemma}[A positive majority expansion for residual compositions]
\label{lem:resglobalmajority}
The function $H(w)=F_D'(\sqrt w)$ is completely monotone on
$(0,\infty)$ and holomorphic on the right half-plane.
It therefore has a representation
\[
H(1-v)=\sum_{k\ge0}d_kv^k,
\qquad d_k\ge0,\qquad \sum_kd_k=\lambda.
\]
If $R=F_D(1)$, then
$F_D(z)=\sum_k(d_k/C_k)M_k(z)$ and
$\sum_kd_k/C_k=R$.
\end{lemma}
\paragraph{Proof idea.}
The derivative of the composed scalar map has a nonnegative expansion in the
majority basis. Closure properties of the squared residual map preserve this
positivity through every layer.

\begin{proof}
Retain $C(w)=\sech^2\sqrt w$ and
$T(w)=\tanh\sqrt w/\sqrt w$ from the scalar proof.
The square of the residual map is
\[
g(w)=p^2w+2p\alpha\sqrt w\tanh(s\sqrt w)
       +\alpha^2\tanh^2(s\sqrt w).
\]
It is Bernstein: each summand is Bernstein, and
\[
\frac d{dw}\bigl[\sqrt w\tanh(s\sqrt w)\bigr]
=\frac s2\{T(s^2w)+C(s^2w)\}
\]
is completely monotone.
Also the one-step derivative
$p+\alpha s C(s^2w)$ is completely monotone.
The closure facts already proved for the scalar composition show
that the product expressing $H$ is completely monotone.

If $\sqrt w=x+iy$ with $\Re w>0$, then $x>|y|$.
Both $x+iy$ and $\tanh(s(x+iy))$ lie in the cone
$\Re z>|\Im z|$: for tanh this follows from
$\sinh(2sx)>|\sin(2sy)|$ and its positive denominator.
Their nonnegative linear combination lies in the same cone.
Thus $g$ maps the right half-plane into itself and the composed
derivative is holomorphic there.
The Taylor expansion and derivative integration from the scalar
case now prove all remaining claims.
\end{proof}

\paragraph{Proof idea.}
Sample one global majority for the whole residual chain and retain its
output-range gate. A rare scan of all inputs supplies the alternative bound
when the input is short.

\begin{proof}[Upper bounds in Theorem~\ref{thm:resmeanlaw}]
Use the global majority mixture and its gate $R$.
At a positive input $z$, the derivative product satisfies
\[
\frac{F_D'(z)}\lambda
=\prod_{j=0}^{D-1}\left(1-\frac{\alpha s}{a}
                       \tanh^2(sF_j(z))\right).
\]
Since $F_j(z)\le a^jz$, the loss from one is at most
$\min\{1,Az^2\}$. The computation of \eqref{eq:chainarity} and
\eqref{eq:chainQ}, with $s^D$ replaced by $\lambda$, bounds the
full-majority source count by $2\lambda\sqrt{1+A}$.
This proves the unknown-source upper bound.

For the actual network output in \eqref{eq:resmeanidentity}, the
unconditional mass of the majority-of-one component is
$d_0=F_D'(1)\ge p^D$.
Replace mass $p^D$ of that component's uniform-coordinate request
by a request for $x_1$. This changes its mean by exactly
$p^D(x_1-M)$ and leaves its cost unchanged.
All mixture masses remain nonnegative.
A second strategy uses the same gate $R$, then scans all $m$
coordinates and evaluates the resulting scalar target.
The modified positive majority representation has total mass
$R$, so it also proves $|h_{D,1}|\le R$ for every context.
This verifies that the gated scan uses a valid conditional sign.
The range lemma gives $mR\le\sqrt5\lambda m/\sqrt{1+A}$.
Select the first strategy when $1+A\le m$, and the second
otherwise. This proves $Q^*\le3\Xi$.
\end{proof}

\paragraph{Proof idea.}
Separate the unknown-source lower bound from the finite-input lower bound.
Endpoint coupling controls the range, while the score of the empirical input
mean captures the derivative growth and its saturation.

\begin{proof}[Lower bounds in Theorem~\ref{thm:resmeanlaw}]
Endpoint coupling gives $U\ge R$ and $Q^*\ge R$.
For the unknown-source problem, if $A<80$, the lower range bound
gives $R\ge\lambda\sqrt{1+A}/81$.
If $A\ge80$, then $T=2R/\lambda\le2/\sqrt{17}<1/2$, and
$F_D'(0)=\lambda$. The bending bound \eqref{eq:bending} with the
scalar cost bound \eqref{eq:unknownlower} yields
\[
U\ge\frac{3\lambda^2}{16R}
\ge\frac3{16\sqrt5}\lambda\sqrt{1+A}
\ge\frac1{81}\lambda\sqrt{1+A}.
\]

For the finite network, under independent input signs of mean $t$,
the expected correction $p^D(X_1-M)$ in
\eqref{eq:resmeanidentity} is zero. Thus its averaged output is
$H(t)=\E_tF_D(M)$.
The weighted Jensen argument used in
\eqref{eq:derivativelower}, applied to the lower range bound, gives
\begin{equation}\label{eq:resfinitehprime}
h:=H'(0)\ge\frac\lambda{\sqrt{1+(3m-2)A/m^2}}
\ge\frac\lambda{\sqrt{1+3A/m}}.
\end{equation}
If $A<400$, endpoint coupling gives $Q^*\ge\Xi/401$.
If $A,m\ge400$, put $T=2R/h$. Then
\[
T^2\le\frac{4(1+3A/m)}{1+A/5}
\le20/A+60/m\le1/5<1/4.
\]
The bending bound \eqref{eq:bending} and the second transcript
bound give
\[
Q^*\ge\frac{3h^2}{16R}
\ge\frac{3\lambda\sqrt{1+A/5}}{16(1+3A/m)}
\ge\frac3{64\sqrt5}\Xi.
\]

It remains to consider $m<400$ and $A\ge400$.
At a balanced input, flipping coordinate $i$ changes the output
mean by absolute amount
\[
F_D(2/m)+2p^D\mathbf1_{\{i=1\}}-2p^D/m.
\]
These amounts are nonnegative because $F_D(z)\ge p^Dz$.
Their sum is $mF_D(2/m)$.
Lemma~\ref{lem:attentioncoupling}, summed over the coordinates, gives
\[
Q^*\ge\frac m2F_D(2/m)
\ge\frac\lambda{\sqrt{1+4A/m^2}}.
\]
If $A\le m^2$, this is at least $\lambda/\sqrt5$, while
$\Xi\le\lambda\sqrt m<20\lambda$.
If $A>m^2$, it is at least
$m\lambda/\sqrt{5A}\ge\Xi/\sqrt5$.
Every resulting constant is at least $1/512$.
\end{proof}

\begin{corollary}[Residual amplification lower bound]
\label{cor:reslower}
Every exact sampler for the residual mean network satisfies
\begin{equation}\label{eq:resfisherlower}
\sup_x\E[Q\mid x]\ge
\frac{\lambda^2}{1+(3m-2)A/m^2}.
\end{equation}
For $1<s\le2$, put $\delta=s-1$. Then, uniformly in
$0\le\alpha\le1$,
\begin{equation}\label{eq:ressuperlower}
\sup_x\E[Q\mid x]\ge
\frac{2\delta}{2\delta+3s^3}
\min\{(1+\alpha\delta)^{2D},m\}.
\end{equation}
At $s=1$, the matching mean-network law becomes
\begin{equation}\label{eq:rescriticalmean}
\frac1{512}\min\!\left\{\sqrt{1+\alpha D},
                   \frac m{\sqrt{1+\alpha D}}\right\}
\le Q^*\le
3\min\!\left\{\sqrt{1+\alpha D},
                   \frac m{\sqrt{1+\alpha D}}\right\}.
\end{equation}
\end{corollary}
\paragraph{Proof idea.}
Insert the finite-input derivative estimate into the transcript inequality.
Elementary bounds on the accumulated residual gain yield the supercritical
and critical consequences.

\begin{proof}
Square \eqref{eq:resfinitehprime} and apply the first transcript
bound. If $\alpha>0$ and $s>1$, then
\[
A=\frac{\alpha s^3}{a}\frac{a^{2D}-1}{a^2-1}
\le\frac{s^3}{2\delta}\lambda^2.
\]
Use $1+(3m-2)A/m^2\le1+3A/m$ and the elementary minimum
inequality from the feed-forward lower bound.
For $\alpha=0$, the target is $x_1$ and the displayed inequality
holds directly. At $s=1$, $a=\lambda=1$ and $A=\alpha D$.
\end{proof}

These bounds identify the residual scaling that permits
polylogarithmic work at polynomial depth.

\begin{corollary}[Residual depth for polylogarithmic work]
\label{cor:residual-loglog-depth}
Fix $s>1$. Consider residual networks \eqref{eq:residualarchitecture}
of width $n$ with row norms at most $s$, biases of absolute value at
most one, and $\alpha$ and all parameters on the grid
$2^{-b}\mathbb Z$, where $20\lceil\log_2n\rceil\le b\le\log^{O(1)}n$
and $D\le n^{O(1)}$. Let $\alpha=\alpha_n$ and $D=D_n$ vary with $n$.
If $\alpha D=O(\log\log n)$, one uniform exact sampler has expected
online bit work $\log^{O(1)}n$ on every such network and context.
Conversely, if some exact sampler has worst-network, worst-context
expected online bit work $\log^{O(1)}n$, then
$\alpha D=O(\log\log n)$.
\end{corollary}
\paragraph{Proof idea.}
The raw residual sampler costs at most an exponential of a fixed
multiple of $\alpha D$. The mean-network lower bound grows
exponentially in $\alpha D$ until it saturates at the width.

\begin{proof}
Here $\bits=\log^{O(1)}n$. For sufficiency,
Corollary~\ref{cor:largeresidualsharpbits} gives expected online bit
work $O(\bits_s^4b_s^D)$ with $b_s^D\le\exp(4\alpha s^2D)$. This is
polylogarithmic when $\alpha D=O(\log\log n)$.

For necessity, choose a fixed dyadic gain
$s_0\in(1,\min\{s,2\}]$, put $\delta=s_0-1$, and let
$m=2^{\lfloor\log_2n\rfloor}\ge n/2$. For large $n$ the weights
$s_0/m$ lie on the grid, so the residual mean network of gain $s_0$
belongs to the class. By \eqref{eq:ressuperlower} and
$\log(1+\alpha\delta)\ge\alpha\delta/2$, every exact sampler makes
at least $c\min\{e^{\delta\alpha D},m\}$ expected probes on some
input, where $c=2\delta/(2\delta+3s_0^3)$. Each probe costs at least
one bit operation. Suppose the work is at most $\log^Kn$ for a fixed
$K$. Since $cm>\log^Kn$ for all sufficiently large $n$, the minimum
is then the exponential term, and $e^{\delta\alpha D}\le
c^{-1}\log^Kn$. Taking logarithms gives $\alpha D=O(\log\log n)$.
\end{proof}
